% Methods restructured 2026-10-06 (EP comments on methods_EP.docx): overview first, one subsection per calculation,
% technical details moved to SI.tex (Section S1). Red = to complete or check.

We computed two kinds of free-energy differences between microstates and combined them in a thermodynamic network.
Deprotonation free energies, which give microscopic \pka{} values, were obtained from DC-DFTB metadynamics in explicit water for the two imidazole nitrogens of histidine and the four nitrogens of the BV core.
Free energies between microstates with the same net charge (tautomers and conformers) were obtained from DC-DFTB umbrella sampling, for the \textit{syn}--\textit{anti} rotation of propionic acid, the intramolecular proton transfer of malonaldehyde, and the proton transfers and ring rotation of BV.
The network gives the microscopic and macroscopic \pka{} values.
For BV, these free energies are also the reference values for constant-pH MD, which gives the population of each microstate as a function of pH; these populations weight the absorption spectra of the individual microstates to give the pH-dependent spectrum.

We distinguish protonation macrostates, protonation or tautomeric microstates, and conformers.
A protonation macrostate groups all species having the same total proton count and net charge, whereas a microstate additionally specifies the location of each titratable proton.
Where useful, we use \emph{protomer} for microstates distinguished specifically by the protonated or deprotonated atomic site.
Conformers share the same proton count, proton placement, and covalent connectivity but differ in molecular geometry.

Histidine comprises three protonation microstates distributed between two protonation macrostates (Figure~\ref{fig:micro}).
The cationic HIP state contains protons on both imidazole nitrogen atoms, whereas the neutral HID and HIE tautomers retain the proton on $N_{\delta}$ and $N_{\epsilon}$, respectively.
Multiple side-chain dihedral conformers were considered within each microstate.
Propionic acid introduces two neutral carboxylic-acid conformers, denoted PRS and PRA according to the orientation of the O--H bond relative to the carbonyl group.
Malonaldehyde (3-hydroxy-2-propenal) is a prototype of intramolecular proton transfer: its enol proton moves between the two oxygen atoms through an internal O--H$\cdots$O hydrogen bond, without leaving the molecule.\cite{baughcumMicrowaveSpectroscopicStudy1981,wangFulldimensionalQuantumCalculations2008}
The two tautomers connected by this transfer, denoted MAA and MAB according to whether the proton sits on oxygen O$_\mathrm{A}$ or O$_\mathrm{B}$, are equivalent by symmetry.

\textcolor{red}{For BV, we considered 11 microstates of the tetrapyrrolic core expected to be relevant over the pH range examined here: seven cationic and four neutral states. These comprise the cationic, fully N-protonated diketo state; four neutral, mono-deprotonated diketo protomers; and six cationic keto--enol tautomers formed by redistribution of an N-bound proton to the carbonyl oxygen of ring A or D. [TO REVISIT once the BV results are final: keto--enol tautomers were not simulated directly. Update this list and the BV network in Figure~\ref{fig:micro} to match what was sampled.]}
\textcolor{green}{\textbf{Response:} The current calculations do not support this 11-state claim. The reactive DC-DFTB work directly sampled PKK and the four mono-deprotonated diketo products, with occasional unconverged keto--enol endpoints; the configured full-chromophore CpHMD calculation currently includes only PPP, BPP, and CPP for the core. Keto--enol states were not included in CpHMD and have no converged populations. The text and Figure~\ref{fig:micro} should therefore be restricted to the states actually sampled, or the keto--enol structures should be shown separately as observed but unquantified products.}
Each BV microstate is identified by a three-letter code.
The first character specifies the nitrogen protonation pattern: P denotes the fully N-protonated core, whereas A, B, C, or D identifies the deprotonated pyrrolic nitrogen.
The second and third characters specify whether terminal rings A and D, respectively, are in the keto (K) or enol (E) form.
Thus, PKK denotes the fully N-protonated diketo state, BKK the B-ring-deprotonated diketo protomer, and BEK a tautomer in which the B-ring nitrogen is deprotonated and ring A is enolized.
The two propionic acid side chains of BV were kept protonated in all microstates because their protonation state has a negligible effect on the calculated absorption properties.\cite{mainasBiliverdinsPropionicChains2024}

\begin{figure}[ht!]
    \centering
    \includegraphics[width=0.9\linewidth]{figures/fig1.png}
    \caption{Definition and connectivity of the microstates considered in this work. The microstate network distinguishes neutral (blue), cationic (teal) and anionic (purple) states; black lines indicate changes in protonation state associated with microscopic \pka{} values, whereas gray lines indicate tautomerization associated with tautomerization constants.}
    \label{fig:micro}
\end{figure}

\subsection{System preparation}

Each solute was placed in a cubic box of TIP3P water\cite{jorgensenComparisonSimplePotential1983} with a buffer distance of 5.5~\AA{} using AmberTools' \texttt{tleap},\cite{case_ambertools_2023} giving 506 atoms for histidine (HIP), \textcolor{red}{XXX for propionic acid (PRS?), XXX for malonaldehyde (MAA),}\textcolor{green}{\textbf{Response:} 428 atoms for propionic acid (PRS) and 372 atoms for malonaldehyde (MAA).} and 869 for BV (PKK, in the locked-helix \textit{ZZZsss} conformation\cite{tongManyRoadsSeam2026}). \textcolor{red}{[Lefteris's BV folders are named solv\_4.0 and the others solv\_5.5: was BV solvated with a 4.0~\AA{} buffer? If so, update this sentence.]}\textcolor{green}{\textbf{Response:} Yes. The saved preparation specification used for these BV runs records a 4.0~\AA{} buffer; histidine, propionic acid, and malonaldehyde used 5.5~\AA{}. The sentence should distinguish BV from the three validation systems.}
The cationic systems (HIP and PKK) were neutralized by a hydroxide ion placed at the solvent position farthest from the solute.
Histidine was described with ff19SB,\cite{tianFf19SBAminoAcidSpecificProtein2020} propionic acid and malonaldehyde with \textcolor{red}{XXX}\textcolor{green}{\textbf{Response:} GAFF2, using the supplied MOL2 charges and system-specific \texttt{frcmod} parameters. The repository does not record the upstream charge-fitting calculation for these two molecules, so that provenance must be recovered before making a more specific charge-method claim.}, and BV with a custom force field derived with ff-opt,\cite{curtisRobustAutomatedForce2026} with RESP\cite{baylyWellbehavedElectrostaticPotential1993} charges fitted to the B3LYP/6-31G* electrostatic potential.
Each system was minimized, heated to 300~K, and equilibrated for 100~ps in the NVT ensemble and 100~ps in the NPT ensemble with Amber\cite{caseAmber20252025} (details in the SI, Section~S1).

The equilibrated structures were the starting points for all quantum-mechanical simulations, in which both the solute and the water were described with the divide-and-conquer formulation\cite{nakai_divide-and-conquer_2023} of third-order self-consistent-charge density-functional tight binding (DC-DFTB3),\cite{gausDensityFunctionalTight2014} as implemented in DCDFTBMD.\cite{nishimura_dcdftbmd_2019}
We used the 3OB Slater--Koster parameter set with the standard X--H damping correction ($\zeta=4.0$),\cite{gaus_parameterization_2014} periodic boundary conditions, and Ewald electrostatics.
Each system was first relaxed at this level for 40~ps at 300~K in the NVT ensemble, with a time step of 1~fs. \textcolor{red}{[Thermostat? The old text said Andersen, 0.01~ps$^{-1}$ (now in the SI); confirm and name it here.]}
\textcolor{green}{\textbf{Response:} Confirmed from the archived DCDFTBMD inputs: \texttt{NVTTYPE=5}, \texttt{BATHTEMP=300}, and \texttt{COLLISIONFREQ=0.01} were used, i.e., the Andersen thermostat with collision frequency 0.01~ps$^{-1}$.}

\subsection{Deprotonation free energies from metadynamics}

Microscopic \pka{} values were obtained from DC-DFTB metadynamics,\cite{barducciMetadynamics2011} following the strategy introduced by Nakai and co-workers for direct \pka{} prediction in explicit water.\cite{Sakti2018}
The collective variable (CV) was the dimensionless coordination number of the titratable nitrogen with all hydrogen atoms,\cite{Tummanapelli2015Amino,Tummanapelli2015Polyprotic}
\begin{equation}
n_{\mathrm{NH}} = \sum_i
\frac{1-\left(r_{\mathrm{N}H_i}/r_0\right)^6}
{1-\left(r_{\mathrm{N}H_i}/r_0\right)^{12}},
\label{eq:coordination-number}
\end{equation}
where the sum runs over all hydrogen atoms $i$ in the system, $r_{\mathrm{N}H_i}$ is the distance between the selected nitrogen and hydrogen atom $i$, and $r_0=1.6$~\AA{} sets the distance over which the coordination function decays.
Because the sum includes every hydrogen atom, biasing $n_{\mathrm{NH}}$ promotes proton loss from the selected nitrogen without prescribing the proton acceptor or the transfer pathway.
We biased each of the two nitrogens of HIP and each of the four nitrogens of PKK, adding Gaussian hills with a height of $5 \times 10^{-4}$~Hartree and a width of 0.1 every 40~fs.
For each site, 200 independent trajectories were started from equilibrated configurations with different initial velocities and random seeds, to sample alternative proton-transfer pathways.
Each trajectory was propagated until $n_{\mathrm{NH}}$ entered the diffusive regime (defined under Trajectory analysis below) and then for an additional 2~ps, giving total trajectory lengths of \textcolor{red}{XXX--YYY}~ps.
\textcolor{green}{\textbf{Response:} This is not how the archived runs were controlled and the sentence should be corrected. The input files specify a fixed upper simulation length; completed runs were subsequently analyzed only through $t_{\mathrm{diff}}+1.75$~ps. The selected handoff trajectories have unequal available lengths (approximately 8.8--20~ps of metadynamics), so there is no defensible universal XXX--YYY range for all 200 trials without first auditing every run.}
Each trajectory required approximately 18.5 computing hours for histidine and 3 days for BV, using 16 cores.
The free-energy profile $F(s)$ along $s=n_{\mathrm{NH}}$ was reconstructed with the standard metadynamics estimator, $F(s,t)\simeq-V_{\mathrm{bias}}(s,t)+C(t)$, and the deprotonation free energy was taken as the difference between the minima of the deprotonated ($s\approx0$) and protonated ($s\approx1$) basins, averaged over ten equally spaced times in the 1.75~ps following the onset of diffusion. \textcolor{red}{[Taken from the old 3$\times$3 summary text; check, including 1.75 vs 2~ps and $T=313.15$~K there vs 300~K here.]}
\textcolor{green}{\textbf{Response:} Confirmed for the final BV and histidine analyses: ten equally spaced samples over the 1.75~ps interval beginning at $t_{\mathrm{diff}}$, with \pka{} converted at 300~K. The old 313.15~K conversion was a legacy analysis setting and has been replaced. One PRN trajectory (syn-6) uses a 1.25~ps window because only that much post-diffusion data was available; it should be disclosed if the PRN metadynamics values remain in the paper.}
Diffusion along $n_{\mathrm{NH}}$ was used only as a sampling and stopping criterion; whether the proton had actually left the solute was assessed separately for each trajectory, as described under Trajectory analysis below.

\subsection{Tautomer and conformer free energies from umbrella sampling}

To evaluate the free-energy differences between tautomers and between conformers, we used umbrella sampling\cite{kastnerUmbrellaSampling2011} with the same DC-DFTB setup as the metadynamics, for consistency.
Intramolecular proton transfer was sampled along a linear combination of distances (LCOD),
\begin{equation}
d_{\mathrm{LCOD}} = d_{\mathrm{H}\cdots\mathrm{D}} - d_{\mathrm{H}\cdots\mathrm{A}},
\label{eq:lcod}
\end{equation}
where $d_{\mathrm{H}\cdots\mathrm{D}}$ is the distance between the transferring proton and the site that holds it in the starting microstate (the donor) and $d_{\mathrm{H}\cdots\mathrm{A}}$ is its distance to the acceptor.
$d_{\mathrm{LCOD}}$ is negative when the proton is bound to the donor, positive when it is bound to the acceptor, and close to zero at the midpoint of the transfer.
This coordinate was used for the MAA$\rightarrow$MAB transfer in malonaldehyde (donor O$_\mathrm{A}$, acceptor O$_\mathrm{B}$) and for the proton transfers between adjacent rings of BV (AKK$\rightarrow$BKK, BKK$\rightarrow$CKK, and CKK$\rightarrow$DKK).
The relative free energies of the four BV protomers were thus obtained in two independent ways, from the metadynamics deprotonation free energies (PKK$\rightarrow$XKK) and from umbrella sampling between adjacent rings, and the two were compared.
Conformational changes were sampled along a dihedral angle: the O=C--O--H dihedral that converts PRS (0$^\circ$) into PRA (180$^\circ$), and the rotation of ring \textcolor{red}{XXX}\textcolor{green}{\textbf{Response:} ring D} about the \textcolor{red}{XXX}\textcolor{green}{\textbf{Response:} C4C--C15 bond, using the NC--C4C--C15--C1D (``single15'' or $\psi_{CD}$) dihedral} that converts \textit{ZZZsss} into \textit{ZZZssa} BV,\cite{tongManyRoadsSeam2026} starting from both PKK and CKK.
The latter gives the deprotonation free energy of the \textit{ZZZssa} conformer through a thermodynamic cycle (see Thermodynamic analysis of the microstate network below).

Windows were spaced by 0.1~\AA{} along $d_{\mathrm{LCOD}}$, from $-1.5$ to $+1.5$~\AA{} for malonaldehyde (31 windows), from $-2.0$ to $+2.0$~\AA{} for the APP$\leftrightarrow$BPP and CPP$\leftrightarrow$DPP transfers of BV (41 windows each), and from $-2.0$ to $+2.3$~\AA{} for the BPP$\leftrightarrow$CPP transfer (44 windows). Dihedral windows were spaced by 10$^\circ$: from 0 to 180$^\circ$ for propionic acid (19 windows) and from $-20$ to 180$^\circ$ for the BV ring-D rotation (21 windows).

The LCOD restraints used the quadratic potential $V=\kappa(\Delta d)^2$ with $\kappa=500$~kJ~mol$^{-1}$~\AA$^{-2}$ (119.50~kcal~mol$^{-1}$~\AA$^{-2}$). Starting structures for the dihedral windows were generated by sequential restrained Amber MD, using 1~ps per window and a force constant of 50~kcal~mol$^{-1}$~rad$^{-2}$. The DC-DFTB dihedral restraints used coefficients of 0.01523087~kJ~mol$^{-1}$~deg$^{-2}$ during equilibration and 0.05~kJ~mol$^{-1}$~deg$^{-2}$ during production, corresponding to 11.95 and 39.23~kcal~mol$^{-1}$~rad$^{-2}$, respectively. Starting structures for the LCOD windows were selected as the frames with $d_{\mathrm{LCOD}}$ closest to each target value from completed DC-DFTB trajectories.

All restrained DC-DFTB simulations were performed at 300~K with a time step of 1~fs. Propionic-acid windows comprised 40~ps of equilibration followed by 80~ps of production, whereas the BV torsional windows comprised 40~ps of equilibration followed by 40~ps of production. The exploratory LCOD profiles were constructed from the available restrained trajectories of nominally 40~ps per window, without a separate production stage. No frames were discarded from the trajectories used for WHAM.

Free-energy profiles were reconstructed using the weighted histogram analysis method (WHAM),\cite{kumarMultidimensionalFreeenergyCalculations1995} and sampling uncertainties were estimated by independently applying circular moving-block bootstrap resampling to each window, using 500 replicas and 2~ps blocks. Overlap between neighboring windows is shown in the Supporting Information. Together, these umbrella-sampling calculations required approximately $2.3\times10^5$ core-hours, based on 16 MPI ranks per window and the recorded elapsed times.

\subsection{Trajectory analysis}
\label{sec:traj-analysis}

Because the CV does not follow the released proton, each metadynamics trajectory was analyzed to identify the location and fate of the released proton.
We followed the excess-proton charge defect through the Mulliken charges of the solute and solvent heteroatoms.
In each frame, the defect was assigned to the solvent oxygen with the largest Mulliken charge within $-0.625\leq q_{\mathrm O}\leq-0.525\ e$, the range that identifies a hydronium-like species.
Correlated charge changes at a donor and an acceptor atom identified intramolecular proton transfer, either direct or through intervening water molecules.
A trajectory was considered to have sampled solute deprotonation only when the charge defect progressed beyond the contact-ion-pair region \textcolor{red}{[defined how? e.g., N--defect distance $>$ XXX~\AA{}]}, migrated away from the solute into the solvent, and did not recombine with either the deprotonated site or the hydroxide counterion.
\textcolor{green}{\textbf{Response:} The implemented operational cutoff is 4.0~\AA{} from the assigned defect oxygen to the nearest solute atom, evaluated with periodic minimum-image distances; it is not specifically an N--defect distance. Passing the cutoff alone is insufficient: endpoint identity and recombination are checked separately.}

The onset of the diffusive regime, $t_{\mathrm{diff}}$, was assigned as the time at which the proton had separated from the nitrogen ($s\leq0.05$ for at least \textcolor{red}{XXX}\textcolor{green}{\textbf{Response:} 50}~fs) and $s$ then recovered to $s\geq0.20$ for a sustained period. \textcolor{red}{[From the old 3$\times$3 summary text; check the persistence interval.]}\textcolor{green}{\textbf{Response:} The automated detector uses 0.05~ps for both persistence tests. However, several final selected BV/PRN $t_{\mathrm{diff}}$ values were manually adjudicated from the full defect diagnostics; the Methods should distinguish the automated candidate criterion from final manual selection.}
To check that diffusion along $s$ after $t_{\mathrm{diff}}$ corresponds to a proton that has moved into the solvent, we also characterized the hydrogen-bonded water wire between the defect and the deprotonated nitrogen: its connectivity (the number of bridging water molecules, or no connection) and its orientation, measured by a conduction-readiness index $\xi$ that equals 1 when all hydrogen bonds along the wire point from the defect toward the nitrogen.
Definitions of these quantities and per-trajectory diagnostic summaries are given in the SI (Section~S1).

\subsection{Thermodynamic analysis of the microstate network}
\label{sec:network}

The sampled free-energy differences were assembled into a network in which each node is a microstate and each edge an equilibrium between two microstates.\cite{ullmannRelationsProtonationConstants2003,bombardaPHDependentPKaValues2010}
Edges between microstates that differ by one proton give microscopic \pka{} values,
\begin{equation}
\mathrm{p}K_{\mathrm{a},\,i\rightarrow j}
=
\frac{\Delta G_{i\rightarrow j}^{\mathrm{deprot}}}{2.303\,RT},
\label{eq:microscopic-pka}
\end{equation}
where $\Delta G_{i\rightarrow j}^{\mathrm{deprot}}$ is the free-energy difference between the protonated basin and the basin containing the deprotonated solute and the solvated excess proton.\cite{Sakti2018}
The free energy of the excess-proton basin is referenced to the aqueous standard state (pH~0) following Tummanapelli and co-workers,\cite{Tummanapelli2014} so that Equation~\ref{eq:microscopic-pka} gives absolute \pka{} values and no experimental \pka{} enters as a calibration.
Edges between microstates with the same net charge give tautomerization or conformational free energies, $\Delta G_{i\rightarrow j}^{\mathrm{taut}}=G_j-G_i$, and the corresponding equilibrium constants $\tau_{i\rightarrow j}=p_j/p_i=\exp(-\Delta G_{i\rightarrow j}^{\mathrm{taut}}/RT)$, where $p_i$ and $p_j$ are the populations of the two microstates.
Because free energy is a state function, the free-energy changes around any closed cycle of the network sum to zero.
Edges sampled in more than one way therefore test the internal consistency of the calculations, and edges that were not sampled follow from cycle closure.

The macroscopic \pka{} combines all microstates of the protonated ($+$) and deprotonated ($0$) manifolds.
Choosing reference microstates $i_0\in+$ and $j_0\in0$,
\begin{equation}
\mathrm{p}K_{\mathrm{a}}^{\mathrm{macro}}
=
\mathrm{p}K_{\mathrm{a},\,i_0\rightarrow j_0}
+
\log_{10}\!\left(\frac{Z_{+}}{Z_{0}}\right),
\qquad
Z_{q}=\sum_{k\in q}\exp\left[-\frac{G_k-G_{k_0}}{RT}\right],
\label{eq:macroscopic-pka}
\end{equation}
where $k_0$ is the reference microstate of manifold $q$.
The macroscopic \pka{} therefore depends not only on the free energy of an individual deprotonation pathway, but also on the number and relative stability of all microstates in each manifold.
The derivation and the general form of these relations are given in the SI (Section~S2).

\subsection{Microstate populations from constant-pH molecular dynamics}
\label{sec:cphmd}

To obtain the population of each BV microstate as a function of pH, which is needed to weight the microstate spectra, we used constant-pH molecular dynamics (CpHMD) in explicit solvent with discrete protonation states,\cite{swailsConstantPHReplica2014,pieriCpHMDThenQMMMIdentification2019} as implemented in Amber,\cite{caseAmber20252025} following the protocol we previously used to titrate the propionic acid side chains of BV.\cite{mainasBiliverdinsPropionicChains2024}
In CpHMD, the system evolves by classical MD, and at regular intervals a Monte Carlo step attempts to change the microstate of the solute, accepting or rejecting the change according to its free-energy cost at the chosen pH.
\textcolor{red}{[Which microstates were included (the PKK and XKK protomers only, or also keto--enol tautomers), and in which conformer (\textit{ZZZsss} only)?]}
\textcolor{green}{\textbf{Response:} No keto--enol tautomers are present in the configured CpHMD state sets. Pairwise capped-core tests include PPP with one of APP/BPP/CPP/DPP; the current full-chromophore BLA CpHMD configuration includes only PPP, BPP, and CPP for the tetrapyrrole core, plus the two titratable propionic tails. The source geometry is the locked-helical starting structure, but the CpHMD input does not restrain a named \textit{ZZZsss} conformer; that assignment should be confirmed from the full-system structure before stating it.}

Setting up CpHMD for BV required three ingredients for each microstate.
First, bonded parameters, taken from the ff-opt force fields described above.\cite{curtisRobustAutomatedForce2026}
Second, atomic point charges.
Atoms whose charge was assumed to change little between microstates were given common (backbone) charges, and only the remaining atoms carried microstate-specific charges.
\textcolor{red}{[Details needed: how the charges were fitted (RESP on which geometries, at which level of theory) and which atoms were treated as backbone.]}
\textcolor{green}{\textbf{Response:} The repository does not contain the original electronic-structure/RESP fitting record, so its geometry and level of theory cannot be recovered here. What is documented is that the capped TPS charge sets were projected from ``biliresp'' PPP/APP/BPP/CPP/DPP MOL2 charges; C3B and C2C were forced to the common PPP junction charge, cap atoms H2C/H3B were set to zero, and the residual needed for an integer state charge was distributed over 12 methyl hydrogens. Thus, the present text about generic ``backbone'' atoms is inaccurate and should be replaced by this explicit projection procedure once the original RESP provenance is supplied.}
Third, a reference free energy for each change of microstate.
We used the microscopic \pka{} values and tautomerization free energies from the DC-DFTB network described above as reference values, and obtained the corresponding reference energies by thermodynamic integration, so that CpHMD reproduces these values for BV in water.
\textcolor{red}{[TI details: number of $\lambda$ windows, length per window, and how the reference energies were extracted.]}
\textcolor{green}{\textbf{Response:} The CpHMD reference energies currently used were calibrated with Amber's \texttt{finddgref.py} workflow, not extracted by thermodynamic integration; the preceding sentence must therefore be corrected. A separate REFEP calculation uses 16 charge-interpolation replicas, 0.2~ns independent equilibration per replica, and 5~ns Hamiltonian replica-exchange production per replica, but that is a validation/comparison calculation and is not the source of the present CpHMD \texttt{DGREF} values.}

After heating and equilibration as described above, CpHMD simulations were run at pH values from \textcolor{red}{XXX}\textcolor{green}{\textbf{Response:} 2.0 for the PPP/BPP and PPP/CPP capped-core tests; 2.5 for PPP/DPP and the current full-chromophore run} to \textcolor{red}{YYY}\textcolor{green}{\textbf{Response:} 6.5 for PPP/BPP and PPP/CPP; 7.0 for PPP/DPP and the full chromophore} in steps of \textcolor{red}{XXX}\textcolor{green}{\textbf{Response:} 0.5 pH unit}, each for \textcolor{red}{XXX}\textcolor{green}{\textbf{Response:} 1.0~ns per replica in the capped-core tests; the current full-chromophore configuration is only 0.10~ns per replica and is not a final production length}~ns.
A change of microstate was attempted every \textcolor{red}{XXX}\textcolor{green}{\textbf{Response:} 200}~fs, and each accepted change was followed by \textcolor{red}{XXX}\textcolor{green}{\textbf{Response:} 0.2}~ps of solvent relaxation.
\textcolor{red}{[Replica exchange between pH values? If so: number of replicas and exchange attempt frequency.]}
\textcolor{green}{\textbf{Response:} Yes: ten pH replicas, with an exchange attempt after every 1000 MD steps = 2~ps.}
From each simulation, we obtained the deprotonated fraction of BV and fitted it to the Hill equation,
\begin{equation}
f_{\mathrm{deprot}}(\mathrm{pH})
=
\frac{1}{1+10^{\,n\left(\mathrm{p}K_{\mathrm{a}}-\mathrm{pH}\right)}},
\label{eq:hill}
\end{equation}
where $n$ is the Hill coefficient, to obtain the titration curve and the macroscopic \pka{}.
The same simulations give the population $p_k(\mathrm{pH})$ of each microstate at each pH, which weights the microstate spectra (see Absorption spectra below).

\subsection{Absorption spectra}
\label{sec:spectra}

An absorption spectrum was computed for each microstate of BV, each described by its own force field, optimized with the same ff-opt procedure as above.\cite{curtisRobustAutomatedForce2026}
To sample solvent arrangements around BV, we ran classical MD for each microstate in a larger water box (15~\AA{} buffer, one chloride ion added to neutralize the cationic states), with BV kept in the locked-helix \textit{ZZZsss} conformation,\cite{tongManyRoadsSeam2026} and selected 199 equally spaced snapshots from XXX~ns of production MD (details in the SI, Section~S1). \textcolor{red}{[Production length to fill in.]}
\textcolor{green}{\textbf{Response:} The spectra-production trajectories and their duration are not present in this repository or in the \texttt{elisa} handoff, so I cannot verify this value from the available files. It must be supplied from the spectroscopy workflow before the placeholder is replaced.}
Each snapshot was the starting point for a 500-fs QM/MM trajectory (0.5~fs time step, 300~K) with TeraChem,\cite{ufimtsevQuantumChemistryGraphical2008,ufimtsevQuantumChemistryGraphical2009,ufimtsevQuantumChemistryGraphical2009a} with BV in the QM region, described at the CAM-B3LYP\cite{yanai_new_2004}-D3\cite{grimme_consistent_2010}/6-31G* level, and TIP3P water in the MM region, cut to a spherical droplet of about 27~\AA{} radius.
This functional was chosen because it has proved to consistently produce accurate bilin geometries.\cite{wiebeler_qmmm_2019}
Excitation energies $\Delta E$ and oscillator strengths $f$ of the lowest 10 singlet excited states were computed on the last frame of each trajectory at the CAM-B3LYP/cc-pVDZ level,\cite{wiebeler_qmmm_2019} and the spectrum of microstate $k$ was obtained by Gaussian broadening,
\begin{equation}
I_k(E)
=
\frac{1}{N}
\sum_{s=1}^{N}
\sum_{n=1}^{10}
f_{k,sn}
\exp\left[
-\frac{\left(E-\Delta E_{k,sn}\right)^2}{2\sigma^2}
\right],
\label{eq:microstate-spectrum}
\end{equation}
where $s$ runs over the $N=199$ sampled geometries, $n$ runs over the excited states, $E$ is the photon energy, and $\sigma=0.10$~eV is the broadening, chosen to match the experimental spectral width of BV in methanol.\cite{liuUltrafastSpectroscopyBiliverdin2020}

The calculated spectrum of an individual microstate is not directly comparable to the experimental spectrum, which reflects the equilibrium ensemble present at a given pH.
The pH-dependent spectrum (calculated between pH~\textcolor{red}{XXX}\textcolor{green}{\textbf{Response:} Not recoverable from the archived spectroscopy data; the configured CpHMD ladders begin at pH 2.0 or 2.5, but that does not establish the spectral grid.} and \textcolor{red}{YYY}\textcolor{green}{\textbf{Response:} Likewise pending the population-weighted spectral dataset.} at intervals of \textcolor{red}{XXX}\textcolor{green}{\textbf{Response:} Likewise pending; do not infer this interval from the 0.5-unit CpHMD replica spacing unless the spectra were actually evaluated only at replica pH values.}) was therefore constructed as the population-weighted sum of the microstate spectra,
\begin{equation}
I_{\mathrm{ens}}(E;\mathrm{pH})
=
\sum_k p_k(\mathrm{pH})\, I_k(E),
\label{eq:ensemble-spectrum}
\end{equation}
where $p_k(\mathrm{pH})$ is the population of microstate $k$ from the CpHMD simulation at that pH.
Spectral maxima were extracted from the population-weighted spectra.
